% ---------------------------------------------------------------- fonts
\usepackage{amsmath,amsthm}
\usepackage{fontspec}
\setmainfont{texgyrepagella}[Extension=.otf,UprightFont=*-regular,
  BoldFont=*-bold,ItalicFont=*-italic,BoldItalicFont=*-bolditalic]
\setmonofont{Inconsolatazi4}[Scale=0.9,Extension=.otf,
  UprightFont=*-Regular,BoldFont=*-Bold]
\usepackage{unicode-math}
\setmathfont{texgyrepagella-math.otf}

% ---------------------------------------------------------------- page
\usepackage[a4paper,top=2.6cm,bottom=2.8cm,inner=3.0cm,outer=3.0cm,headheight=14pt]{geometry}
\usepackage{microtype}
\usepackage{xcolor}
\usepackage{graphicx}
\usepackage{booktabs,array,longtable,tabularx,multirow}
\usepackage[shortlabels]{enumitem}
\usepackage{fancyhdr}
\usepackage{titlesec}
\usepackage{fancyvrb}
\usepackage{caption}
\captionsetup{font=small,labelfont=bf,margin=1em}
\usepackage[most]{tcolorbox}
\usepackage{tikz}
\usetikzlibrary{positioning,arrows.meta,calc,decorations.pathreplacing,fit,shapes.misc}
\usepackage[hidelinks,bookmarksnumbered]{hyperref}
\setlist{itemsep=2pt,topsep=4pt}
\setlength{\parskip}{3pt plus 1pt}
\linespread{1.04}

\definecolor{codegray}{gray}{0.40}
\definecolor{rulegray}{gray}{0.55}
\definecolor{shade}{gray}{0.88}
\definecolor{shadedark}{gray}{0.72}

% ---------------------------------------------------------------- headers
\pagestyle{fancy}
\fancyhf{}
\fancyhead[C]{\small Data Structures \textbar{} 2nd Long Exam Reviewer}
\fancyfoot[C]{\small Juevesano D. 2026 All rights reserved.}
\fancyfoot[LE,RO]{\small\thepage}
\renewcommand{\headrulewidth}{0.4pt}
\fancypagestyle{plain}{\fancyhf{}%
  \fancyhead[C]{\small Data Structures \textbar{} 2nd Long Exam Reviewer}%
  \fancyfoot[C]{\small Juevesano D. 2026 All rights reserved.}%
  \fancyfoot[LE,RO]{\small\thepage}%
  \renewcommand{\headrulewidth}{0.4pt}}

% ---------------------------------------------------------------- titles
\titleformat{\chapter}[display]
  {\centering\normalfont}
  {\large\textsc{\chaptertitlename\ \thechapter}}
  {10pt}
  {\LARGE\scshape}
  [\vspace{6pt}\titlerule]
\titlespacing*{\chapter}{0pt}{10pt}{28pt}
\titleformat{\section}{\normalfont\large\bfseries}{\thesection}{0.8em}{}
\titleformat{\subsection}{\normalfont\normalsize\bfseries\itshape}{\thesubsection}{0.8em}{}
\titleformat{\part}[display]{\centering\normalfont}{\Large\textsc{\partname\ \thepart}}{14pt}
  {\Huge\scshape}[\vspace{10pt}\rule{0.3\linewidth}{0.4pt}]
\renewcommand{\thepart}{\Roman{part}}
\newcommand{\para}[1]{\par\smallskip\noindent\textbf{\textit{#1.}}\enspace}

% ---------------------------------------------------------------- theorems
\theoremstyle{definition}
\newtheorem{definition}{Definition}[chapter]
\newtheorem{example}[definition]{Example}
\theoremstyle{plain}
\newtheorem{theorem}[definition]{Theorem}
\newtheorem{proposition}[definition]{Proposition}
\newtheorem{corollary}[definition]{Corollary}
\theoremstyle{remark}
\newtheorem{remark}[definition]{Remark}

% A note with a bar on the left: \begin{note}{Title} ... \end{note}
\newtcolorbox{note}[1]{enhanced,breakable,frame hidden,boxrule=0pt,
  borderline west={1.2pt}{0pt}{black},colback=white,left=8pt,right=0pt,
  top=1pt,bottom=1pt,before upper={\textsc{#1.}\enspace},fontupper=\small,
  before skip=8pt,after skip=8pt}

% ---------------------------------------------------------------- code
\usepackage{listings}
\lstdefinestyle{book}{
  language=C,
  basicstyle=\ttfamily\small,
  keywordstyle=\bfseries,
  commentstyle=\color{codegray},
  stringstyle=\color{black},
  columns=fullflexible,keepspaces=true,
  numbers=left,numberstyle=\tiny\color{codegray},numbersep=8pt,
  frame=tb,framerule=0.4pt,rulecolor=\color{black},
  xleftmargin=16pt,framexleftmargin=0pt,
  aboveskip=8pt,belowskip=6pt,
  breaklines=true,breakatwhitespace=true,
  showstringspaces=false,tabsize=4,
  morekeywords={bool,true,false,class,public,private,new,delete,NULL,
                using,namespace,typedef,struct},
  upquote=true,
}
\lstset{style=book}
% Boxed starter code in problem statements
\lstdefinestyle{starter}{style=book,numbers=none,frame=single,
  xleftmargin=0pt,framexleftmargin=4pt,framexrightmargin=4pt,
  basicstyle=\ttfamily\footnotesize,aboveskip=6pt,belowskip=4pt}
\newcommand{\codepath}{../code}
% \listingfile[options]{path}{caption}: a library or solution file (shown whole)
\newcommand{\listingfile}[3][]{\lstinputlisting[#1,caption={#3}]{\codepath/#2}}
% \programfile[options]{path}: a program whose first line (// libs:) is hidden
\newcommand{\programfile}[2][]{\lstinputlisting[#1,firstline=2]{\codepath/#2}}
% \outputfile{name}: the real output, out/<name>.txt
\newcommand{\outputfile}[1]{%
  \VerbatimInput[frame=single,framesep=6pt,fontsize=\small,rulecolor=\color{rulegray}]{\codepath/out/#1.txt}}
\renewcommand{\lstlistingname}{Listing}
\newcommand{\cc}[1]{\texttt{#1}}

% ---------------------------------------------------------------- exam items
\makeatletter
\newcommand{\difficulty}[1]{\hfill\textit{\mdseries #1}}
% \problem{n}{Title}{Easy}
\newcommand{\problem}[3]{\par\bigskip\needspace{8\baselineskip}%
  \phantomsection\def\@currentlabel{#1}\label{prob:#1}%
  \addcontentsline{toc}{subsection}{Problem #1. #2}%
  \noindent{\large\bfseries Problem #1.}\enspace\textsc{#2}\difficulty{#3}\par
  \vspace{-4pt}\noindent\rule{\linewidth}{0.4pt}\par\vspace{2pt}}
% \tracing{n}{Title}
\newcommand{\tracing}[2]{\par\bigskip\needspace{8\baselineskip}%
  \phantomsection\def\@currentlabel{#1}\label{trace:#1}%
  \addcontentsline{toc}{subsection}{Tracing #1. #2}%
  \noindent{\large\bfseries Tracing #1.}\enspace\textsc{#2}\par
  \vspace{-4pt}\noindent\rule{\linewidth}{0.4pt}\par\vspace{2pt}}
% \solution{Problem}{n}
\newcommand{\solution}[2]{\par\bigskip\needspace{8\baselineskip}%
  \phantomsection\addcontentsline{toc}{subsection}{Solution to #1 #2}%
  \noindent{\large\bfseries Solution to #1 #2.}\par
  \vspace{-4pt}\noindent\rule{\linewidth}{0.4pt}\par\vspace{2pt}}
\newcommand{\field}[1]{\par\noindent\textsc{#1.}\enspace}
\newcommand{\permitted}{\field{Permitted operations}}
\newcommand{\constraints}{\field{Constraints}}
\newcommand{\required}{\field{Required}}
\newcommand{\bonus}{\field{Bonus (extra credit)}}
\newcommand{\idea}{\field{Idea}}
\newcommand{\complexity}{\field{Complexity}}
\newcommand{\mistakes}{\field{Common mistakes}}
\makeatother
% sample tests: \begin{tests} input & output \\ ... \end{tests}
\newenvironment{tests}
  {\par\smallskip\noindent\small
   \tabularx{\linewidth}{@{}>{\ttfamily\raggedright\arraybackslash}X>{\ttfamily\raggedright\arraybackslash}X@{}}
   \toprule \normalfont\textsc{Input} & \normalfont\textsc{Expected output}\\ \midrule}
  {\bottomrule\endtabularx\par\smallskip}
\usepackage{needspace}
% concept questions
\newlist{questions}{enumerate}{1}
\setlist[questions]{label=\textbf{\arabic*.},ref=\arabic*,leftmargin=*,itemsep=6pt}

% ---------------------------------------------------------------- tikz styles
\tikzset{
  cell/.style={draw,minimum width=7mm,minimum height=6.5mm,inner sep=0pt,font=\small},
  sh/.style={fill=shade},
  shd/.style={fill=shadedark},
  nd/.style={draw,minimum width=6mm,minimum height=6mm,inner sep=1pt,font=\small},
  ptr/.style={-{Stealth[length=2mm]},thin},
  lbl/.style={font=\ttfamily\scriptsize},
  idx/.style={font=\tiny,text=codegray},
}
% \llnode{name}{at}{item}: a singly linked node drawn as [item|.]
\newcommand{\llnode}[3]{%
  \node[nd] (#1) at (#2) {#3};
  \node[draw,minimum width=3mm,minimum height=6mm,inner sep=0pt,anchor=west] (#1p) at (#1.east) {};}

\setlength{\emergencystretch}{3em}
